\documentclass[11pt]{article}

\usepackage[landscape,margin=0.8in]{geometry}

\usepackage{amsmath,amssymb}

\usepackage{graphicx}

\usepackage{booktabs}

\usepackage{caption}

\usepackage{float}

\usepackage{xcolor}

\usepackage[hidelinks]{hyperref}

\usepackage{etoolbox}

\usepackage{longtable}

% Start every top-level section on a new page

\pretocmd{\section}{\clearpage}{}{}

% Allows variable names with underscores without escaping

\newcommand{\code}[1]{\texttt{\detokenize{#1}}}

\newcommand{\kt}{\kappa_t}

\newcommand{\Nbar}{\bar N}

\newcommand{\Nhat}{\hat N}

\graphicspath{{../results/figures/}}

\title{\textbf{Market Concentration and the Flattening of the Phillips Curve}\\[2pt]

\large A Theory-Motivated, Semi-Structural State-Space Assessment of the HSA Mechanism}

\author{Satoshi Ichikawa}

\date{Draft --- August 2026}

\begin{document}

\section{Data}

The estimation dataset is constructed in STATA directly from the raw source

files. Each source is first converted into a quarterly dataset indexed by the

STATA quarterly date, after which the component datasets are merged onto a

contiguous quarterly time spine.

Python is not used for the general data build. The only exception is the PCHIP

interpolation of the annual competition series, since STATA does not provide a

native monotone-cubic PCHIP routine. All other transformations, filters, merges,

lags, and final data exports are carried out in STATA.



\subsection{Inflation}

\paragraph{Source.}

Headline CPI, PPI, core CPI, core PCE, and headline PCE are obtained from

\textbf{FRED}.

\paragraph{Series.}

All inflation measures are constructed from monthly price indexes. Monthly

observations are first averaged within calendar quarter.

\paragraph{Construction.}

For each quarterly price index $P_t$, year-on-year inflation is defined as the

four-quarter log difference,

\begin{equation}
\pi_t^{\mathrm{yoy}}
=
100
\left(
\log P_t-\log P_{t-4}
\right).
\end{equation}

The corresponding quarter-on-quarter inflation rate is

\begin{equation}
\pi_t^{\mathrm{qoq}}
=
100
\left(
\log P_t-\log P_{t-1}
\right).
\end{equation}

The same construction is applied to headline CPI, PPI, core CPI, core PCE, and

headline PCE. For quarterly NKPC specifications, the quarter-on-quarter measure

is the natural counterpart to the one-period inflation rate in the model.

Annualized quarter-on-quarter inflation can equivalently be expressed as

\begin{equation}
\pi_t^{\mathrm{qoq,ann}}
=
400
\left(
\log P_t-\log P_{t-1}
\right).
\end{equation}



\subsection{Inflation Expectations}

We retain inflation expectations at their documented forecast horizons. The
one-year-ahead series are used as broader short-run expectation measures, while
the SPF one-quarter-ahead GDP price-index forecast provides the expectation
measure that maps most directly into a quarterly NKPC.

\subsubsection{Cleveland Fed Inflation Expectations}

\paragraph{Source.}
Inflation expectations are obtained from the Cleveland Fed inflation-expectation
series.

\paragraph{Series.}
The monthly one-year-ahead expectation series is converted to percentage units
and averaged within calendar quarter.

\paragraph{Construction.}
The Cleveland Fed expectation is retained at its original one-year horizon. We
do not divide the one-year-ahead rate by four to construct a one-quarter-ahead
forecast, because such a transformation would not in general recover
$E_t\pi_{t+1}$.

\subsubsection{Survey of Professional Forecasters}

\paragraph{Source.}
Inflation expectations are also obtained from the
\textbf{Survey of Professional Forecasters (SPF)} of the Federal Reserve Bank
of Philadelphia.

\paragraph{One-year-ahead forecasts.}
The SPF one-year-ahead file reports GDP-price-index and CPI inflation
expectations. These are median-based expectations for average inflation over
the four quarters following the quarter in which the survey is conducted. The
GDP-price-index series is constructed from the median level forecasts for the
survey quarter and four quarters ahead, while the CPI series is constructed
from the quarterly CPI inflation forecasts over the following four quarters.
The published one-year-ahead series are retained directly, with no frequency
conversion or interpolation.

\paragraph{One-quarter-ahead GDP price-index forecast.}
For the quarterly NKPC, we additionally construct the SPF one-quarter-ahead GDP
price-index inflation forecast from the \emph{mean} level forecasts. In the SPF
level-file nomenclature, PGDP2 denotes the mean forecast of the GDP price-index
level for the quarter in which the survey is conducted, while PGDP3 denotes the
mean forecast for the following quarter. The SPF official mean growth rate is
the growth rate of these mean level forecasts; it is not the cross-sectional
mean of individual growth-rate forecasts.

Following the SPF quarterly-growth convention, the one-quarter-ahead forecast
is therefore

\begin{equation}
\widehat{\pi}_{t+1|t}^{\mathrm{SPF},1q}
=
100\left[
\left(
\frac{\widehat{P}_{t+1|t}^{\mathrm{mean}}}
     {\widehat{P}_{t|t}^{\mathrm{mean}}}
\right)^4-1
\right],
\end{equation}

where the SPF expresses quarterly growth in annualized percentage points using
discrete compounding.

Because the NKPC inflation measures in this dataset are also retained in
annualized log-difference form, we additionally construct the log-equivalent
expectation

\begin{equation}
\widehat{\pi}_{t+1|t}^{\mathrm{SPF},1q,\log}
=
400\left(
\log \widehat{P}_{t+1|t}^{\mathrm{mean}}
-
\log \widehat{P}_{t|t}^{\mathrm{mean}}
\right).
\end{equation}

The first measure follows the SPF official quarterly-growth convention; the
second is a transformation into the continuous-compounding units used by the
annualized log-inflation specifications and is not itself an official SPF
reported growth series.

\subsection{Output Gap}

We consider several alternative measures of real marginal cost and economic

slack that can enter the Phillips curve forcing term.



\subsubsection{Beveridge--Nelson Output Gap}

\paragraph{Source.}

The Beveridge--Nelson output measure is constructed from quarterly real GDP.

\paragraph{Series.}

The decomposition provides a transformed output series together with its

cyclical component.

\paragraph{Construction.}

The corresponding trend is reconstructed as

\begin{equation}
y_t^{\mathrm{BN,trend}}
=
y_t^{\mathrm{BN}}
-
y_t^{\mathrm{BN,gap}}.
\end{equation}



\subsubsection{Hodrick--Prescott Output Gap}

\paragraph{Source.}

The HP output gap is constructed from quarterly real GDP.

\paragraph{Series.}

Real output is transformed into 100 log points before filtering.

\paragraph{Construction.}

Let $Y_t$ denote real output. The transformed series is

\begin{equation}
y_t
=
100\log Y_t.
\end{equation}

The quarterly HP filter uses the conventional smoothing parameter

\begin{equation}
\lambda=1600.
\end{equation}

The decomposition is

\begin{equation}
y_t
=
\tau_t^{\mathrm{HP}}
+
c_t^{\mathrm{HP}},
\end{equation}

where $\tau_t^{\mathrm{HP}}$ denotes the trend and

$c_t^{\mathrm{HP}}$ denotes the cyclical output gap.



\subsubsection{Inverse Markup}

\paragraph{Source.}

The business markup series is taken from Nekarda and Ramey.

\paragraph{Series.}

Quarterly observations are available directly from the source. Because the

markup exhibits substantial low-frequency variation, we also use a

Beveridge--Nelson decomposition to isolate its cyclical component.

\paragraph{Construction.}

The inverse markup is defined as

\begin{equation}
x_t
=
\frac{1}{\mu_t},
\end{equation}

where $\mu_t$ denotes the markup.

The cyclical inverse-markup measure is obtained from the Beveridge--Nelson

decomposition and provides an alternative empirical proxy for real marginal

cost.



\subsubsection{Labor-Share Gap}

\paragraph{Source.}

The labor-share measure is obtained from the seasonally adjusted labor-share

index available from \textbf{FRED}.

\paragraph{Series.}

The quarterly labor-share index is restricted to positive observations.

\paragraph{Construction.}

Let $L_t$ denote the labor-share index. The transformed series is

\begin{equation}
\ell_t
=
100\log L_t.
\end{equation}

The quarterly HP filter uses

\begin{equation}
\lambda=1600,
\end{equation}

and decomposes the series as

\begin{equation}
\ell_t
=
\tau_{L,t}^{\mathrm{HP}}
+
c_{L,t}^{\mathrm{HP}},
\end{equation}

where $\tau_{L,t}^{\mathrm{HP}}$ denotes the labor-share trend and

$c_{L,t}^{\mathrm{HP}}$ denotes the cyclical labor-share gap.



\subsubsection{Negative Unemployment Gap}

\paragraph{Source.}

The natural rate of unemployment and the seasonally adjusted unemployment rate

are obtained from \textbf{FRED}.

\paragraph{Series.}

Both monthly series are averaged within calendar quarter.

\paragraph{Construction.}

The negative unemployment gap is defined as

\begin{equation}
u_t^{\mathrm{gap}}
=
u_t^{*}-u_t,
\end{equation}

where $u_t^{*}$ denotes the natural rate of unemployment and $u_t$ denotes the

observed unemployment rate.

A positive value therefore corresponds to unemployment being below its estimated

natural rate.



\subsection{Competition Measures}

We use three measures of product-market competition constructed from distinct

data sources. The first is based on the industry-concentration data of Grullon,

Larkin, and Michaely (2019). The second is based on the Hoberg--Phillips

Text-based Network Industry Classifications (TNIC). The third is constructed

directly from quarterly firm-level revenue data from S\&P Capital IQ.



\subsubsection{Grullon--Larkin--Michaely Measure}

\paragraph{Source.}

The first competition measure is based on the industry-concentration data

associated with Grullon, Larkin, and Michaely (2019),

\textit{Are US Industries Becoming More Concentrated?}

\paragraph{Series.}

The original measure is observed at annual frequency. We retain the level series

together with its Beveridge--Nelson cycle and trend components.

\paragraph{Construction.}

For quarterly-frequency specifications, the annual series and its

Beveridge--Nelson components are interpolated to quarterly frequency using

monotone-cubic PCHIP interpolation.

For mixed-frequency specifications, the annual observations are instead used

without interpolation. The observed annual value is assigned to Q4 of the

corresponding year, while Q1--Q3 are left missing.

The two representations therefore serve different purposes. The interpolated

series provide a quarterly path for quarterly-frequency estimation, while the

Q4-only series preserve the original annual information at its native frequency

for mixed-frequency estimation.

Because the PCHIP knots are located at year-end, the interpolated and

observed-only representations coincide at Q4.



\subsubsection{Hoberg--Phillips TNIC Measure}

\paragraph{Source.}

The second competition measure is based on the Hoberg--Phillips Text-based

Network Industry Classifications (TNIC). TNIC identifies product-market

competitors using pairwise similarity in firms' 10-K product descriptions,

allowing competitor sets to vary across firms.

\paragraph{Series.}

The resulting competition measure is observed at annual frequency. We retain

the level series together with its Beveridge--Nelson cycle and trend components.

\paragraph{Construction.}

For quarterly-frequency specifications, the annual series and its

Beveridge--Nelson components are interpolated to quarterly frequency using

monotone-cubic PCHIP interpolation.

For mixed-frequency specifications, no interpolation is applied. The original

annual observations are retained at Q4, while Q1--Q3 are left missing.

As with the Grullon--Larkin--Michaely measure, the interpolated series is used

for quarterly-frequency estimation and the observed-only series preserves the

native annual observations for mixed-frequency estimation. The two

representations coincide at Q4.



\subsubsection{Capital IQ Competition Measure}

\paragraph{Source.}

The third competition measure is constructed directly from S\&P Capital IQ.

The sample consists of U.S. public firms listed on major exchanges, with

firm-level total revenue observed at quarterly frequency.

\paragraph{Series.}

Unlike the two annual competition measures, the Capital IQ measure is

constructed natively at quarterly frequency and therefore requires no temporal

interpolation.

Firm-quarter observations with missing reported revenue are excluded. Negative

reported revenues are floored at zero when constructing market shares, while

the original reported values are retained separately.

Industry classification is based on SIC. A firm's current SIC classification is

applied to its full revenue history.

\paragraph{Market definition.}

The baseline broad-market definition is based on SIC divisions, with

manufacturing separated into durable and nondurable goods so that the entire

manufacturing sector is not treated as a single market.

Alternative market definitions based on SIC, NAICS, and S\&P sector

classifications are also retained in the underlying competition panel.

\paragraph{Firm revenue shares.}

Within market $j$ and quarter $t$, firm $i$'s revenue share is

\begin{equation}
s_{ijt}
=
\frac{R_{ijt}}
{\sum_{i\in j}R_{ijt}},
\end{equation}

where $R_{ijt}$ denotes cleaned quarterly revenue.

\paragraph{Market concentration.}

Market-level concentration is measured by the Herfindahl--Hirschman index,

\begin{equation}
HHI_{jt}
=
\sum_{i\in j}s_{ijt}^{2}.
\end{equation}

The baseline construction retains only market-quarter observations containing

at least 10 firms.

\paragraph{Seasonal adjustment.}

Each market-level HHI is seasonally adjusted before economy-wide aggregation.

The adjustment removes systematic within-year movements in concentration within

each market while retaining idiosyncratic quarter-specific variation.

Let $HHI_{jt}^{\mathrm{SA}}$ denote the resulting seasonally adjusted market

concentration measure.

\paragraph{Economy-wide effective number of firms.}

The economy-wide effective number of firms is defined as the inverse of the

weighted mean market HHI,

\begin{equation}
N_t
=
\frac{1}
{\sum_j w_{jt}HHI_{jt}^{\mathrm{SA}}},
\qquad
\sum_j w_{jt}=1.
\end{equation}

We construct both firm-count-weighted and revenue-weighted versions,

\begin{equation}
N_t^{\mathrm{firm}}
=
\frac{1}
{\sum_j w_{jt}^{\mathrm{firm}}HHI_{jt}^{\mathrm{SA}}},
\end{equation}

and

\begin{equation}
N_t^{\mathrm{rev}}
=
\frac{1}
{\sum_j w_{jt}^{\mathrm{rev}}HHI_{jt}^{\mathrm{SA}}}.
\end{equation}

Because these measures are constructed directly from quarterly firm-level

revenues, no temporal interpolation is required.



\subsection{Establishment Stock}

\paragraph{Source.}

The establishment-stock measure combines the annual Business Dynamics Statistics

(BDS) establishment count with quarterly Business Employment Dynamics (BED)

establishment births and deaths.

\paragraph{Series.}

The BDS data provide a level benchmark for the number of establishments, while

the BED data provide quarterly establishment birth and death flows. The BED

series are reported in thousands and are converted to establishment counts

before constructing the stock.

\paragraph{Construction.}

The 1993 BDS establishment count is used as the level anchor. Quarterly

establishment net entry is defined as

\begin{equation}
\Delta E_t
=
B_t-D_t,
\end{equation}

where $B_t$ denotes establishment births and $D_t$ denotes establishment deaths.

The establishment stock is then constructed as

\begin{equation}
E_t
=
E_{1993}^{\mathrm{BDS}}
+
\sum_{s\leq t}\Delta E_s.
\end{equation}



\subsection{Lags and Timing}

For regressors entering the empirical specifications with a one-quarter lag, we

also retain the corresponding lagged observation. For a generic quarterly

series $z_t$,

\begin{equation}
z_t^{\mathrm{prev}}
=
z_{t-1}.
\end{equation}

Quarterly observations are dated at the end of the corresponding calendar

quarter.



\subsection{Variable Definitions}

Table~\ref{tab:data_variables} summarizes the variable names used in the

model-ready dataset. Variable names are reported here, rather than in the

preceding discussion, to separate the economic definitions from their

implementation in the data files.

\small

\begin{longtable}{p{0.16\textwidth}p{0.27\textwidth}p{0.51\textwidth}}

\caption{Variables in the Model-Ready Dataset}

\label{tab:data_variables}\\

\toprule

Category & Variable & Description \\

\midrule

\endfirsthead

\multicolumn{3}{c}{\tablename\ \thetable\ -- continued from previous page} \\

\toprule

Category & Variable & Description \\

\midrule

\endhead

\midrule

\multicolumn{3}{r}{Continued on next page} \\

\endfoot

\bottomrule

\endlastfoot

Inflation

& \code{pi_cpi}

& Headline CPI inflation, year-on-year log difference. \\

Inflation

& \code{pi_cpi_qoq}

& Headline CPI inflation, quarter-on-quarter log difference. \\

Inflation

& \code{pi_cpi_qoq_ann}

& Annualized quarter-on-quarter headline CPI log inflation. \\

Inflation

& \code{pi_cpi_core}

& Core CPI inflation, year-on-year log difference. \\

Inflation

& \code{pi_cpi_core_qoq}

& Core CPI inflation, quarter-on-quarter log difference. \\

Inflation

& \code{pi_cpi_core_qoq_ann}

& Annualized quarter-on-quarter core CPI log inflation. \\

Inflation

& \code{pi_pce}

& Headline PCE inflation, year-on-year log difference. \\

Inflation

& \code{pi_pce_qoq}

& Headline PCE inflation, quarter-on-quarter log difference. \\

Inflation

& \code{pi_pce_qoq_ann}

& Annualized quarter-on-quarter headline PCE log inflation. \\

Inflation

& \code{pi_pce_core}

& Core PCE inflation, year-on-year log difference. \\

Inflation

& \code{pi_pce_core_qoq}

& Core PCE inflation, quarter-on-quarter log difference. \\

Inflation

& \code{pi_pce_core_qoq_ann}

& Annualized quarter-on-quarter core PCE log inflation. \\

Inflation

& \code{pi_ppi}

& PPI inflation, year-on-year log difference. \\

Inflation

& \code{pi_ppi_qoq}

& PPI inflation, quarter-on-quarter log difference. \\

Inflation

& \code{pi_ppi_qoq_ann}

& Annualized quarter-on-quarter PPI log inflation. \\

\midrule

Expectations

& \code{Epi}

& Cleveland Fed one-year-ahead inflation expectation, retained at its original horizon. \\

Expectations

& \code{Epi_spf_gdp}

& SPF one-year-ahead GDP-price inflation expectation; median-based average over the four quarters following the survey quarter. \\

Expectations

& \code{Epi_spf_cpi}

& SPF one-year-ahead CPI inflation expectation; median-based average over the four quarters following the survey quarter. \\

Expectations

& \code{Epi_spf_gdp_1q}

& SPF one-quarter-ahead GDP price-index growth constructed from mean level forecasts using the official annualized discrete Q/Q convention. \\

Expectations

& \code{Epi_spf_gdp_1q_log}

& Annualized log-change equivalent of the SPF one-quarter-ahead mean-level forecast; not an official SPF reported growth series. \\

\midrule

Competition

& \code{N_Gustavo}

& PCHIP-interpolated quarterly Grullon--Larkin--Michaely competition series. \\

Competition

& \code{N_Gustavo_BN_cycle}

& Beveridge--Nelson cycle of the Grullon--Larkin--Michaely series, interpolated quarterly. \\

Competition

& \code{N_Gustavo_BN_trend}

& Beveridge--Nelson trend of the Grullon--Larkin--Michaely series, interpolated quarterly. \\

Competition

& \code{N_Gustavo_annual_q4}

& Original annual Grullon--Larkin--Michaely observation retained at Q4 only. \\

Competition

& \code{N_Gustavo_BN_cycle_annual_q4}

& Original annual Beveridge--Nelson cycle retained at Q4 only. \\

Competition

& \code{N_Gustavo_BN_trend_annual_q4}

& Original annual Beveridge--Nelson trend retained at Q4 only. \\

Competition

& \code{N_TNIC}

& PCHIP-interpolated quarterly Hoberg--Phillips TNIC competition series. \\

Competition

& \code{N_TNIC_BN_cycle}

& Beveridge--Nelson cycle of the TNIC series, interpolated quarterly. \\

Competition

& \code{N_TNIC_BN_trend}

& Beveridge--Nelson trend of the TNIC series, interpolated quarterly. \\

Competition

& \code{N_TNIC_annual_q4}

& Original annual TNIC observation retained at Q4 only. \\

Competition

& \code{N_TNIC_BN_cycle_annual_q4}

& Original annual TNIC Beveridge--Nelson cycle retained at Q4 only. \\

Competition

& \code{N_TNIC_BN_trend_annual_q4}

& Original annual TNIC Beveridge--Nelson trend retained at Q4 only. \\

Competition

& \code{N_capitaliq_firmw}

& Capital IQ effective number of firms using firm-count-weighted market HHIs. \\

Competition

& \code{N_capitaliq_revw}

& Capital IQ effective number of firms using revenue-weighted market HHIs. \\

\midrule

Output gap

& \code{output_BN}

& Beveridge--Nelson transformed output series. \\

Output gap

& \code{output_gap_BN}

& Beveridge--Nelson cyclical output gap. \\

Output gap

& \code{output_trend_BN}

& Beveridge--Nelson output trend. \\

Output gap

& \code{output}

& Log real output. \\

Output gap

& \code{output_trend_HP}

& Hodrick--Prescott output trend. \\

Output gap

& \code{output_gap_HP}

& Hodrick--Prescott cyclical output gap. \\

Output gap

& \code{labor_share}

& Quarterly labor-share index. \\

Output gap

& \code{labor_share_100log}

& Labor-share index expressed in 100 log points. \\

Output gap

& \code{labor_share_trend_HP}

& Hodrick--Prescott labor-share trend. \\

Output gap

& \code{labor_share_gap_HP}

& Hodrick--Prescott labor-share gap. \\

Output gap

& \code{unemp_gap}

& Negative unemployment gap, natural rate minus observed unemployment. \\

Output gap

& \code{markup}

& Nekarda--Ramey business markup. \\

Output gap

& \code{markup_inv}

& Inverse markup. \\

Output gap

& \code{markup_BN_inv}

& Beveridge--Nelson cyclical inverse-markup measure. \\

\midrule

Establishments

& \code{establishment_births}

& Quarterly establishment births. \\

Establishments

& \code{establishment_deaths}

& Quarterly establishment deaths. \\

Establishments

& \code{establishment_net_entry}

& Quarterly establishment births minus deaths. \\

Establishments

& \code{establishment_stock}

& Establishment stock constructed from the 1993 BDS anchor and cumulative BED net entry. \\

\midrule

Lag

& \code{pi_cpi_prev}

& One-quarter lag of headline CPI inflation. \\

Lag

& \code{pi_ppi_prev}

& One-quarter lag of PPI inflation. \\

Lag

& \code{pi_cpi_core_prev}

& One-quarter lag of core CPI inflation. \\

Lag

& \code{pi_pce_prev}

& One-quarter lag of headline PCE inflation. \\

Lag

& \code{pi_pce_core_prev}

& One-quarter lag of core PCE inflation. \\

Lag

& \code{pi_cpi_qoq_prev}

& One-quarter lag of quarter-on-quarter headline CPI inflation. \\

Lag

& \code{pi_ppi_qoq_prev}

& One-quarter lag of quarter-on-quarter PPI inflation. \\

Lag

& \code{pi_cpi_core_qoq_prev}

& One-quarter lag of quarter-on-quarter core CPI inflation. \\

Lag

& \code{pi_pce_qoq_prev}

& One-quarter lag of quarter-on-quarter headline PCE inflation. \\

Lag

& \code{pi_pce_core_qoq_prev}

& One-quarter lag of quarter-on-quarter core PCE inflation. \\

Lag

& \code{pi_cpi_qoq_ann_prev}

& One-quarter lag of annualized quarter-on-quarter headline CPI log inflation. \\

Lag

& \code{pi_ppi_qoq_ann_prev}

& One-quarter lag of annualized quarter-on-quarter PPI log inflation. \\

Lag

& \code{pi_cpi_core_qoq_ann_prev}

& One-quarter lag of annualized quarter-on-quarter core CPI log inflation. \\

Lag

& \code{pi_pce_qoq_ann_prev}

& One-quarter lag of annualized quarter-on-quarter headline PCE log inflation. \\

Lag

& \code{pi_pce_core_qoq_ann_prev}

& One-quarter lag of annualized quarter-on-quarter core PCE log inflation. \\

Lag

& \code{unemp_gap_prev}

& One-quarter lag of the negative unemployment gap. \\

Lag

& \code{markup_inv_prev}

& One-quarter lag of the inverse markup. \\

Lag

& \code{markup_BN_inv_prev}

& One-quarter lag of the Beveridge--Nelson cyclical inverse markup. \\

Lag

& \code{output_gap_BN_prev}

& One-quarter lag of the Beveridge--Nelson output gap. \\

Lag

& \code{output_gap_HP_prev}

& One-quarter lag of the Hodrick--Prescott output gap. \\

Lag

& \code{labor_share_gap_HP_prev}

& One-quarter lag of the Hodrick--Prescott labor-share gap. \\

\end{longtable}

\normalsize

\section{Summary of Data Construction}

The main data choices can be summarized as follows:

\begin{itemize}

    \item \textbf{Data construction.}

    The model-ready dataset is built in STATA directly from the raw source

    files. The only exception is the PCHIP interpolation of the annual

    competition measures, which is performed outside STATA.

    \item \textbf{Inflation.}

    Inflation is constructed from quarterly averages of monthly price indexes.

    Year-on-year, quarter-on-quarter, and annualized quarter-on-quarter log

    inflation measures are retained. The quarterly measures correspond most

    directly to the one-period inflation rate in the quarterly NKPC, while the

    year-on-year measure is retained as an alternative specification.

    \item \textbf{Inflation expectations.}

    Expectations are retained at their documented horizons rather than being
    mechanically converted across horizons.

    \begin{itemize}

        \item Cleveland Fed expectations are retained at the one-year-ahead
        horizon; no division by four is used to approximate a next-quarter
        expectation.

        \item SPF one-year-ahead GDP-price and CPI expectations are the
        published median-based measures of average inflation over the four
        quarters following the survey quarter.

        \item The SPF one-quarter-ahead GDP price-index measure is constructed
        from PGDP2 and PGDP3 in the mean level forecasts. Consistent with the
        SPF documentation, the official-style measure is the annualized
        discrete Q/Q growth rate of the mean forecast levels, not the mean of
        individual growth-rate forecasts.

        \item A separate annualized log-change equivalent of the one-quarter-
        ahead SPF forecast is retained for specifications using annualized log
        inflation.

    \end{itemize}

    \item \textbf{Phillips-curve forcing variable.}

    Several alternative measures of economic slack or real marginal cost are

    retained:

    \begin{itemize}

        \item Beveridge--Nelson output gap;

        \item Hodrick--Prescott output gap;

        \item inverse markup and its Beveridge--Nelson cyclical component;

        \item Hodrick--Prescott labor-share gap; and

        \item negative unemployment gap.

    \end{itemize}

    This allows the Phillips-curve results to be evaluated without relying on a

    single empirical proxy for the theoretical forcing variable.

    \item \textbf{Grullon--Larkin--Michaely competition measure.}

    The original measure is annual. For quarterly-frequency specifications, the

    annual level and its Beveridge--Nelson components are interpolated using

    monotone-cubic PCHIP. For mixed-frequency specifications, the original annual

    observations are retained at Q4 without interpolation.

    \item \textbf{Hoberg--Phillips TNIC competition measure.}

    The TNIC-based measure is also annual and is treated in the same way:

    PCHIP interpolation is used for quarterly-frequency specifications, while

    the original Q4 observations are used for mixed-frequency specifications.

    \item \textbf{Capital IQ competition measure.}

    The Capital IQ measure is constructed directly at quarterly frequency from

    revenues of U.S. public firms listed on major exchanges. Firm-level revenue

    shares are aggregated into market-level HHIs, which are seasonally adjusted

    before aggregation. Economy-wide effective-firm measures are then constructed

    using both firm-count and revenue weights.

    \item \textbf{Establishment dynamics.}

    Quarterly establishment births and deaths from BED are combined with the

    1993 BDS establishment count as a level anchor. The resulting establishment

    stock provides a measure of changes in the population of active

    establishments that is distinct from the revenue-based concentration

    measures.

    \item \textbf{Timing and lags.}

    Quarterly observations are dated at calendar-quarter end. One-quarter lags

    are retained for variables entering lagged specifications.

\end{itemize}

\end{document}
